\documentclass[11pt, a4paper]{article}

% ── 字体与编码 ─────────────────────────────────────────────
\usepackage{fontspec}
\usepackage[fontset=fandol]{ctex}
\setmainfont{Times New Roman}

% ── 页面布局 ──────────────────────────────────────────────
\usepackage[a4paper, margin=2.0cm, top=2.2cm, bottom=2.2cm]{geometry}
\usepackage{setspace}
\setstretch{1.2}

% ── 数学 ──────────────────────────────────────────────────
\usepackage{amsmath}

% ── 表格 ──────────────────────────────────────────────────
\usepackage{booktabs}
\usepackage{float}
\usepackage[table]{xcolor}
\usepackage{array}

% ── 其他 ──────────────────────────────────────────────────
\usepackage{enumitem}
\usepackage{titlesec}
\usepackage{fancyhdr}
\usepackage{tcolorbox}
\tcbuselibrary{breakable, skins}
\usepackage[colorlinks=true, linkcolor=blue!60!black, urlcolor=blue!60!black]{hyperref}
\usepackage{tikz}
\usetikzlibrary{arrows.meta, shapes.geometric, positioning, fit}

% ── 标题样式 ──────────────────────────────────────────────
\titleformat{\section}{\large\bfseries\heiti}{\thesection.}{0.5em}{}[\vspace{-0.3em}]
\titleformat{\subsection}{\normalsize\bfseries\heiti}{\thesubsection}{0.5em}{}
\titlespacing{\section}{0pt}{1em}{0.4em}
\titlespacing{\subsection}{0pt}{0.6em}{0.2em}

% ── 颜色框 ────────────────────────────────────────────────
\newtcolorbox{notebox}[1][]{
  colback=gray!6, colframe=gray!50,
  title=#1, fonttitle=\bfseries\small\heiti,
  boxrule=0.5pt, left=4pt, right=4pt, top=3pt, bottom=3pt,
  breakable
}

% ─────────────────────────────────────────────────────────
\begin{document}

% ── 总标题页 ────────────────────────────────────────────────
\begin{center}
  {\LARGE\bfseries\heiti 广州花园酒店评论分析技术报告}\\[0.5em]
  {\normalsize }
\end{center}
\vspace{0.8em}

% ╔══════════════════════════════════════════════════════════
% ║  第一章 评论聚类分析
% ╚══════════════════════════════════════════════════════════
\begin{center}
  {\Large\bfseries\heiti 第一章\quad 评论聚类分析}\\[0.3em]
  {\small}
\end{center}
\vspace{0.5em}

% ─────────────────────────────────────────────────────────
\section{背景与目标}
% ─────────────────────────────────────────────────────────

广州花园酒店的评论数据已有人工标注的14个类别标签（如"餐饮设施"、"前台服务"、"地铁交通"等），平均每条评论对应2.80个标签，其中97.7\%的评论具有2个及以上标签。人工标注虽然准确，但存在固有局限：标注成本高、类别边界主观、难以动态扩展，且标注规则并不总能反映评论文本的实际聚集模式。本部分的目标是用数据驱动的无监督聚类方法重新发现评论主题，验证与人工分类的一致性，并解决多标签分配问题，最终输出带结构化标签的聚类结果。

% ─────────────────────────────────────────────────────────
\section{技术方案}
% ─────────────────────────────────────────────────────────

\subsection{文本向量化：TF-IDF}

中文评论首先经jieba分词，过滤停用词（含通用停用词及酒店专有无意义词"花园酒店"、"入住"等），再用TF-IDF（\texttt{max\_features=3000, min\_df=3, max\_df=0.85, sublinear\_tf=True}）转化为3000维稀疏向量。选择TF-IDF而非词嵌入（如Word2Vec）的原因在于：\textbf{聚类不需要跨语义泛化，TF-IDF对主题词的区分更直接}；同时，TF-IDF的稀疏矩阵便于后续c-TF-IDF判别词提取。

\subsection{降维：UMAP}

由于高维空间中欧氏距离趋于均匀，密度概念失效，因此3000维TF-IDF向量不能直接聚类。选用UMAP而非PCA是因为UMAP能保留\textbf{局部邻域结构}（非线性流形），而PCA只保留全局方差方向，对文本主题的局部聚集性捕捉较差。

UMAP分两阶段进行：第一阶段降至\textbf{50维}（\texttt{min\_dist=0.0}），使同类评论更紧密，供HDBSCAN聚类；第二阶段降至\textbf{2维}（\texttt{min\_dist=0.1}），适当扩散，供人眼可视化。两次UMAP均用余弦距离（\texttt{metric='cosine'}），因为文本向量的方向比绝对长度更能反映语义相似性。

\subsection{聚类：HDBSCAN}

选择HDBSCAN而非K-Means的核心理由有三：(1) \textbf{无需预设簇数}$k$，让数据自行决定有多少个主题；(2) \textbf{支持噪声点}，那些内容太综合或太边缘的评论不会被强行归类，而是标记为噪声（本数据集中占44.1\%）；(3) \textbf{适应非球形簇}，评论主题的分布在语义空间中并非规则圆形。最终发现14个聚类，与人工标注类别数巧合一致，但这一结果并非预设。

参数选择：\texttt{min\_cluster\_size=30}控制最小簇规模（低于30条的话题视为噪声），\texttt{min\_samples=5}控制核心点敏感度，\texttt{cluster\_selection\_method='eom'}（超额质量法）倾向于生成语义更纯的、数量适中的簇。

% ─────────────────────────────────────────────────────────
\section{关键工程问题与解决方案}
% ─────────────────────────────────────────────────────────

\subsection{问题一：施工噪音簇误匹配餐饮内容}

初始多标签分配发现，大量含餐饮内容的语义段被错误分配到施工噪音簇（簇8）。根本原因是：簇8的TF-IDF均值中心包含"早餐"、"自助"、"晚餐"等词，普通均值中心无法区分"这是噪音簇偶然包含的食物词"还是"食物词是噪音簇的本质特征"。

解决方案是引入\textbf{c-TF-IDF（簇间IDF加权）}：
$$
\text{c-TF-IDF}_{w,c} = \text{TF-IDF}_{w,c} \times \log\!\left(1 + \frac{N_c}{df_w + 1}\right)
$$
其中$df_w$为词$w$出现在几个簇中。"早餐"在多个簇都有权重，$df_w$大，簇间IDF小，乘下来权重被压低；而"施工"、"太吵"只在簇8出现，$df_w$小，权重被放大。每簇取top-30判别词构成稀疏判别中心后，簇8的代表词变为：施工、噪音、太吵、不好、成就，食物词完全被淘汰。此外还加入后处理规则：若某段落的最佳匹配是施工噪音簇，但段落中不含"施工/噪音/太吵"等实词，则强制降级到次优匹配。

\subsection{问题二：餐饮类别缺失}

HDBSCAN未产生独立的"餐饮"簇。原因是含餐饮内容的评论往往同时提及其他话题（住宿体验、地理位置等），在语义空间中较为分散，密度不足以聚集成簇。然而餐饮是业务上最重要的类别之一。

解决方案是引入\textbf{手工餐饮伪簇}（ID=100）：从TF-IDF词表中直接构造一个餐饮中心向量，包含早餐、自助、晚餐、餐厅、菜品、食材等餐饮核心词，同时将行政/酒廊等词清零，避免与行政酒廊簇（簇2）抢竞争。这是一种半监督补充：用已知业务知识弥补纯数据驱动的盲区，最终236条评论（10.9\%）被分配到餐饮标签。

\subsection{多标签分配机制}

为解决一条评论涉及多个主题的问题，将每条评论用DashScope嵌入（\texttt{text-embedding-v3}，1024维）按话题跳变处切分为若干语义段（平均2.2段/条），每段独立与判别中心矩阵做余弦匹配，汇总取频次最高的最多3个标签。最终平均1.49个标签/条，40.5\%的评论获得$\geq$2个标签。

% ─────────────────────────────────────────────────────────
\section{LLM语义标签生成}
% ─────────────────────────────────────────────────────────

HDBSCAN只输出簇编号，需要将每个簇转化为业务可读的标签。我们采用DashScope \texttt{qwen-turbo}，向LLM提供每簇的top-20判别词和3条代表性评论样例，要求生成"【一级分类 - 二级细分】"格式的标签（一级固定为设施/服务/位置/体验四类）。

为防止LLM对相似簇生成重复标签，每次调用都将已生成的所有标签传入prompt。即便如此，LLM仍产生了4处错误（如将前台服务簇（簇3）标注为"房间升级"，将整体口碑簇（簇5）标注为已存在的"前台服务"重复标签），均通过Patch单元格直接覆盖修正。最终15个标签（14个HDBSCAN簇+1个手工餐饮簇）语义清晰、无重复。

% ─────────────────────────────────────────────────────────
\section{与人工标签的对比分析}
% ─────────────────────────────────────────────────────────

\subsection{类别对应关系}

聚类结果与人工标注在\textbf{主要类别上高度吻合}，但两者在粒度和覆盖上存在差异（见表~\ref{tab:cluster-compare}）。

\begin{table}[H]
\centering
\caption{聚类标签与人工标签对应关系}
\label{tab:cluster-compare}
\setlength{\tabcolsep}{5pt}
\begin{tabular}{lll}
\toprule
\textbf{聚类标签} & \textbf{对应人工标签} & \textbf{备注} \\
\midrule
【位置类 - 地铁便利性】（133条） & 地铁交通 & 高度吻合 \\
【设施类 - 餐饮设施】（236条） & 餐饮设施 & 需手工伪簇补充 \\
【服务类 - 前台服务】（59条） & 前台服务 & 高度吻合 \\
【设施类 - 卫浴设施】（96条） & 房间设施 & 聚类更细分（专注卫浴） \\
【设施类 - 行政酒廊】（152条） & 餐饮设施/会员权益 & 人工未单独区分 \\
【体验类 - 施工噪音】（62条） & 整体体验/无单独标签 & 聚类独立发现负面主题 \\
【体验类 - 节日氛围】（67条） & 整体体验/无单独标签 & 聚类独立发现节庆主题 \\
【体验类 - 回头意向】（192条） & 整体满意度 & 聚类捕捉行为信号 \\
【服务类 - 生日客情】（43条）  & 前台服务/特殊需求 & 聚类细粒度更高 \\
\bottomrule
\end{tabular}
\end{table}

\subsection{差异分析}

\textbf{聚类的优势}体现在以下方面。第一，粒度更细：人工标签中"餐饮设施"是一个大类，聚类则将行政酒廊（会员早餐+下午茶权益，152条）和普通餐饮（236条）明确分开，这对于不同业务线的分析更有价值。第二，发现人工未设定的类别：施工噪音（62条）和节日氛围（67条）在人工标注体系中没有独立类别，但聚类算法从数据中自然发现这两个有意义的主题，反映了一批客人的共同体验。第三，行为性标签：回头意向簇（192条，关键词为"下次还会来"等）是一种\textbf{行为意图信号}，在原始人工标签中归为笼统的"整体满意度"，聚类后可单独提取用于预测客户忠诚度。

\textbf{聚类的不足}体现在：平均每条评论仅1.49个标签，远低于人工的2.80个，原因是语义分段匹配的相似度阈值（0.05）过于保守，部分段落未能触发多标签；同时44.1\%的评论被HDBSCAN归为噪声，需通过分段匹配兜底，覆盖率虽然完整，但主簇硬标签缺失。此外，聚类对短评（少于5词后分词）效果较差，这类评论更多依赖兜底规则而非真正的语义匹配。

\subsection{结论}

两套标签体系互为补充：人工标签精度高、覆盖全，适合作为评估基准；聚类标签能发现人工体系盲区、细化粒度、提取行为信号，尤其是能针对酒店特性生成多样化的标签（包含行政酒廊、生日客情等特殊细分模块），适合作为探索性分析和自动化扩展的工具。由于人工标签具备更高的通用性，更匹配C端用户选择酒店时的需求，我们\textbf{在问答系统中仍然采用人工标签}；但在实际应用中，可将聚类标签与人工标签同时保留，供不同下游任务选用。

% ─────────────────────────────────────────────────────────
\section{聚类结果汇总}
% ─────────────────────────────────────────────────────────

\begin{notebox}[聚类分析最终结果]
最终得到\textbf{15个语义聚类标签}（14个HDBSCAN自动发现 + 1个手工餐饮伪簇），全部标签语义清晰、无重复。多标签分配结果：平均1.49标签/条，$\geq$2标签占40.5\%，$\geq$3标签占8.8\%，含餐饮标签236条（10.9\%）。输出文件：\texttt{filtered\_comments\_clustered.csv}（含6个新增列）和 \texttt{category\_summaries\_cluster.json}（15条类别元信息），可直接供RAG问答系统检索使用。
\end{notebox}


\vspace{1.5em}

% ╔══════════════════════════════════════════════════════════
% ║  第二章 评论摘要生成与检索
% ╚══════════════════════════════════════════════════════════
\setcounter{section}{0}
\renewcommand{\thesection}{2.\arabic{section}}
\begin{center}
  {\Large\bfseries\heiti 第二章\quad 评论摘要生成与检索}\\[0.3em]
  {\small }
\end{center}
\vspace{0.5em}

% ─────────────────────────────────────────────────────────
\section{背景与目标}
% ─────────────────────────────────────────────────────────

酒店评论RAG（检索增强生成）问答系统包含五路混合检索机制：BM25稀疏检索、向量相似度检索、反向Query检索、HyDE假设文档扩展，以及\textbf{类别摘要检索}。其中，前四路均直接召回具体评论，而类别摘要检索的目标是：在回答用户问题时，提供超越单条评论的\textbf{宏观类别背景}，帮助LLM做出更全面、准确的综合判断。

本章聚焦于类别摘要的两个核心环节：（1）摘要的离线生成流程——如何从2171条原始评论提炼出14个类别的高质量摘要；（2）摘要的在线检索机制——如何在用户提问时精准召回相关类别摘要，并注入LLM生成答案。

% ─────────────────────────────────────────────────────────
\section{摘要生成}
% ─────────────────────────────────────────────────────────

\subsection{知识库结构与类别划分}

评论数据共2171条，经人工多标签标注覆盖14个业务类别，包括：整体满意度（1855条）、前台服务（1352条）、餐饮设施（705条）、房间设施（619条）、交通便利性（465条）、公共设施（277条）、客房服务（261条）、卫生状况（403条）、周边配套（250条）、退房/入住效率（105条）、性价比（156条）、安静程度（147条）、景观/朝向（88条）、价格合理性（13条）。平均每条评论对应2.80个标签，97.7\%的评论具有2个及以上标签。

类别间样本量差异显著（最多1855条，最少13条），这对摘要质量和向量检索可信度均有影响，是后续优化中需要特别处理的维度。

\subsection{分层聚合生成方法}

直接将类别下全部评论拼接后输入LLM会超出上下文窗口（尤其是1855条整体满意度评论），且大量重复性内容会稀释关键信息。本章采用\textbf{分层聚合}策略（见表~\ref{tab:gen}）：

\begin{enumerate}[leftmargin=2em]
  \item \textbf{分批摘要（Batch Summarization）}：将每类别的评论按批次（每批约30-50条）输入 \texttt{qwen-turbo}，生成批次小结，提取该批次的核心主题与高频词；
  \item \textbf{关键词整合}：对批次小结做关键词频次统计，提取全类别 top-5 代表性关键词（如"前台服务"类的\textit{热情周到、主动升级、员工点名表扬、服务两极化、高峰期效率}）；
  \item \textbf{最终摘要合成（Meta-Summarization）}：将所有批次小结和关键词集中，再次调用LLM生成最终摘要（约500-1000字），要求覆盖正负面评价、指出核心矛盾、保留具体细节。
\end{enumerate}

\begin{table}[H]
\centering
\caption{分层摘要生成流程}
\label{tab:gen}
\setlength{\tabcolsep}{6pt}
\begin{tabular}{clll}
\toprule
\textbf{阶段} & \textbf{输入} & \textbf{模型} & \textbf{输出} \\
\midrule
批次摘要 & 30-50条原始评论 & qwen-turbo & 批次小结 + 关键词 \\
关键词聚合 & 所有批次关键词 & 词频统计 & top-5类别关键词 \\
最终合成 & 全部批次小结 & qwen-turbo & 500-1000字类别摘要 \\
\bottomrule
\end{tabular}
\end{table}

\subsection{摘要质量评价}

最终14条摘要平均长度约650字，覆盖正面评价、负面评价与具体细节三个维度。以"前台服务"类摘要为例，其明确指出服务热情度的两极分化——多位员工被高频实名表扬，同时高峰期效率问题与个别员工短板亦有所记录，体现了评论数据的真实分布。

% ─────────────────────────────────────────────────────────
\section{摘要检索机制}
% ─────────────────────────────────────────────────────────

\subsection{向量化方案：非对称检索}

摘要知识库的构建与查询均基于 \texttt{text-embedding-v4}（DashScope，1024维），但两阶段使用不同的向量化策略，构成\textbf{非对称检索}（Asymmetric Retrieval）：

\begin{itemize}[leftmargin=2em]
  \item \textbf{离线构建（Document Side）}：使用 \texttt{embed\_batch()}，对拼接文本 \texttt{"类别名：关键词。摘要前300字"} 进行普通嵌入，向量充分表达文档语义；
  \item \textbf{在线检索（Query Side）}：使用 \texttt{embed\_query()}，额外传入 \texttt{text\_type="query"} 和任务指令
    \begin{center}
    \textit{"Given a hotel review query, retrieve relevant hotel reviews that answer the query."}
    \end{center}
    该指令激活 \texttt{text-embedding-v4} 的指令微调能力，将查询向量拉向"文档回答"的语义空间，而非"关键词匹配"空间。
\end{itemize}

非对称检索相较于对称检索（query和document用同一策略），在短查询/长文档场景下可显著提升召回率，尤其对"游泳池干净吗"这类口语化问题与长篇摘要文本的匹配效果更为明显。

此外，embedding文本由原始方案的\textbf{仅关键词}（约30字）升级为\textbf{类别名+关键词+摘要前300字}（约330字），使向量包含足够的语义内容，避免了短文本向量退化为关键词集合的问题。

\subsection{优化A：相似度阈值过滤}

原方案固定返回 top-3 摘要，不论余弦相似度高低。在14个类别中，即使最相关的类别相似度仅为0.20，仍会被召回并注入LLM prompt，引入噪声。

优化方案引入 \textbf{相似度阈值}（$\tau = 0.25$）：相似度低于阈值的类别摘要不参与最终结果，候选池同时扩大至8条以覆盖过滤后仍能保留足够结果。此外，若同一类别存在多个匹配文档（如原始摘要与QA变体），按类别去重，只保留最高分的一条。

$$
\text{摘要}_{i} \in \text{结果} \iff \text{cosine\_sim}(q, d_i) \geq \tau
$$

其中 $q$ 为查询向量，$d_i$ 为第 $i$ 个摘要文档向量，$\tau=0.25$ 为过滤阈值。

\subsection{优化B：意图规则混合召回}

纯向量检索对口语化短查询（如"前台怎么样"）存在Gap——由于query向量维度高，在14类摘要构成的小索引中，语义对齐可能不稳定。本项目已有意图检测模块（\texttt{detect\_intent()}），基于关键词词典对查询进行意图分类，但原始实现中摘要检索完全未利用该信号。

优化方案在向量检索的基础上，叠加一层\textbf{意图$\rightarrow$类别规则映射}（表~\ref{tab:rule}）：若意图关键词命中且对应类别未出现在向量结果中，则直接从ChromaDB按类别名查询并强制补入结果，标记来源为 \texttt{summary\_rule}。

\begin{table}[H]
\centering
\caption{意图$\rightarrow$摘要类别映射规则}
\label{tab:rule}
\setlength{\tabcolsep}{6pt}
\begin{tabular}{ll}
\toprule
\textbf{识别到的意图} & \textbf{强制补充的摘要类别} \\
\midrule
早餐 & 餐饮设施 \\
房间 & 房间设施、卫生状况 \\
服务 & 前台服务、客房服务、退房/入住效率 \\
设施 & 公共设施 \\
位置 & 交通便利性、周边配套 \\
\bottomrule
\end{tabular}
\end{table}

最终排序策略：向量召回结果按相似度降序排列，规则补充结果附后，整体截取前3条。这确保了向量检索的精度主导权，同时规则提供兜底保障。

\subsection{优化C：检索导向的QA变体生成}

原始摘要为叙述式长文本（约650字），其向量与用户的短查询之间存在\textbf{Query-Document Gap}：用户提问"游泳池水温怎么样"，但摘要文本以"综合X条评论，酒店公共设施整体表现卓越……"开头，两者的语义空间并不对齐。

优化方案在摘要知识库中为每个类别额外生成\textbf{3条QA变体句}，每句20-50字，风格类似对用户问句的直接回答，分别覆盖：主要优点、需注意的问题、一个具体特色细节。QA变体由 \texttt{qwen-turbo} 根据完整摘要生成，embedding后存入同一ChromaDB集合，\texttt{document} 字段仍指向完整摘要，\texttt{metadata} 中额外保留 \texttt{qa\_text} 字段记录原始变体句。

这样，向量检索命中QA变体时，LLM获得的仍是完整的类别背景摘要，但召回路径更贴近用户问句形式，有效缩短了Query-Document Gap（如图~\ref{fig:qa}所示）。

\begin{figure}[H]
\centering
\begin{tikzpicture}[
  node distance=0.7cm and 1.4cm,
  box/.style={rectangle, draw=gray!70, fill=gray!8, rounded corners=3pt,
              minimum width=3.2cm, minimum height=0.7cm, align=center, font=\small},
  arrow/.style={-Stealth, thick, gray!60}
]
\node[box] (q)   {用户查询\\（短句）};
\node[box, right=1.6cm of q]  (qa)  {QA变体\\（20-50字回答句）};
\node[box, right=1.6cm of qa] (sum) {完整类别摘要\\（650字）};

\draw[arrow] (q)  -- node[above, font=\footnotesize]{向量相似度高} (qa);
\draw[arrow] (qa) -- node[above, font=\footnotesize]{document字段} (sum);
\draw[arrow, dashed, bend right=25] (q) to node[below, font=\footnotesize, gray]{原始Gap大} (sum);
\end{tikzpicture}
\caption{QA变体作为检索桥接层，缩短用户查询与摘要文档之间的语义距离}
\label{fig:qa}
\end{figure}

% ─────────────────────────────────────────────────────────
\section{摘要在LLM生成中的使用}
% ─────────────────────────────────────────────────────────

召回的类别摘要以结构化方式注入LLM prompt，位于具体评论列表之前，形成两层上下文：

\begin{notebox}[Prompt结构]
\textbf{【类别背景知识】}（基于大量评论的综合分析）：\\
【前台服务】基于1352条评论，前台服务是酒店体验的核心亮点……（截取前300字）\\
【房间设施】酒店房间设施整体呈现显著的"新旧并存"特征……（截取前300字）\\[0.5em]
\textbf{【真实住客评论】}：\\
评论1【大床房，评分4.5】: 前台小姐姐超级贴心……\\
评论2【标准房，评分3.0】: 房间有点旧，设施一般……
\end{notebox}

这种"宏观背景 + 具体证据"的双层结构，使LLM能在生成答案时兼顾全局趋势（来自摘要）与个体案例（来自评论），回答质量相比纯评论上下文有明显提升，尤其对"整体上评价怎么样"类泛化问题效果更显著。每条摘要截取前300字注入，避免因摘要过长挤占评论的上下文空间。

% ─────────────────────────────────────────────────────────
\section{摘要结果汇总与改进方向}
% ─────────────────────────────────────────────────────────

\begin{notebox}[摘要检索当前实现总览]
\textbf{知识库}：14条类别摘要（原始）+ 42条QA变体，存储于ChromaDB本地持久化库（\texttt{data/chroma\_db/}），向量维度1024，余弦相似度空间。\\[0.4em]
\textbf{检索机制}：非对称向量检索（候选池8条）$\rightarrow$ 余弦相似度阈值过滤（$\tau=0.25$）$\rightarrow$ 类别去重 $\rightarrow$ 意图规则补充 $\rightarrow$ 截取top-3。\\[0.4em]
\textbf{生成注入}：召回摘要以"类别背景知识"块前置注入prompt，每条限300字，与真实评论构成双层上下文。
\end{notebox}

\vspace{0.6em}

\textbf{主要改进对比}（见表~\ref{tab:summary-compare}）：

\begin{table}[H]
\centering
\caption{摘要检索机制改进前后对比}
\label{tab:summary-compare}
\setlength{\tabcolsep}{5pt}
\begin{tabular}{p{3.5cm}p{5.0cm}p{5.5cm}}
\toprule
\textbf{维度} & \textbf{改进前} & \textbf{改进后} \\
\midrule
检索方式 & \texttt{get()}全量拉取，固定返回前3条 & \texttt{query()}向量相似度检索，8候选池 \\
Embedding文本 & 仅关键词（约30字） & 类别名+关键词+摘要前300字（约330字） \\
查询向量化 & \texttt{embed\_batch()}无指令 & \texttt{embed\_query()}加task指令（非对称） \\
相关性过滤 & 无，始终返回3条 & 余弦相似度阈值$\tau=0.25$过滤 \\
意图利用 & 未使用意图信号 & 规则映射强制补充命中类别 \\
LLM上下文 & 摘要从未传入LLM & 双层prompt（背景+评论），300字截断 \\
QA变体 & 无 & 每类别3条变体句，缩短Query Gap \\
\bottomrule
\end{tabular}
\end{table}

\textbf{后续可优化方向}：

\begin{enumerate}[leftmargin=2em]
  \item \textbf{阈值自适应}：当前阈值$\tau=0.25$为固定值，可根据意图置信度动态调整——高置信度查询收紧阈值以保证精准，低置信度放宽以增加覆盖；
  \item \textbf{评论数量加权}：对余弦相似度叠加评论数量权重
    $\text{score}' = \text{score} \times (1 + \alpha \log(\text{comment\_count}))$，
    避免13条评论的"价格合理性"类别与1352条评论的"前台服务"类别在同等条件下获得相同权重；
  \item \textbf{扩展Query投票}：利用已有的 \texttt{expand\_intent()} 对多个扩展问题分别检索，对类别分数加权平均，进一步提升召回稳定性。
\end{enumerate}

\end{document}
